%!TEX root = ../hutbthesis_main.tex

\chapter{绪论}

\section{研究背景与意义}

虚拟现实、机器人控制、远程医疗、虚拟培训等领域中，可变形物体（如软组织、柔性材料、布料、器官等）的交互仿真的技术越来越成为关键技术依赖的对象，刚体物体的建模和交互技术比较成熟，但是可变形物体受到力的作用时会表现出非线性力学特性，产生较大的变形，其参数随时间变化，接触状态也较为复杂，因此要实现高精度、高时效性的建模与交互仍然存在很大困难，柔性体的力触觉建模复杂度远高于刚体模型，传统建模方式难以适配动态形变的交互需求，无法满足精细仿真作业要求。\cite{ref32}仅仅依靠视觉信息很难支持精细的操作，因为操作人员缺乏力觉和触觉的感受，无法准确判断接触力、变形程度和表面纹路，导致操作不稳、任务完成率低，严重制约了远程手术、精密装配、康复训练等场景的实际应用\cite{ref31}。

在此背景下，将依靠物理引擎的实时仿真与触觉反馈技术相结合，这是冲破前面瓶颈的关键途径，MuJoCo凭借高效多体动力学求解能力、优秀的柔性体仿真性能，在可变形物体建模、接触计算、实时性保障方面具有明显优势，能有效适配柔性交互场景的动态仿真需求，目前已广泛应用于机器人柔顺交互、精密操作仿真等相关研究领域\cite{ref33}。WEART触觉手套可以给指尖提供压力、纹理、温度等多模态的触觉反馈，它们互相配合就形成了一个视觉、物理仿真、触觉反馈的闭环系统。

本课题以MuJoCo可变形物体实时触摸建模与仿真系统为研究对象，围绕可变形物体触摸交互中实时形变建模、接触力精确计算、触觉信号映射、视触同步、量化评价等主要问题展开研究，试图构建出一个高刷新率、低时延且可以进行定量评定的视触融合仿真平台，从而为可变形物体交互原理的研究、算法验证和应用系统开发提供标准化、可重现的实验环境，具有重要的理论价值和工程意义。

\section{国内外研究现状}

\subsection{国内研究现状}

国内学者对于可变形物体建模、触觉渲染、XR交互、医疗机器人、生物力学、复合材料仿真等各个方面已经形成了比较系统全面的研究架构，包含理论方法、硬件设计到系统实现的一整套技术。

胡子阳等人提出了一种结合CBAM注意力机制的人体软组织多物理场快速同步建模方法，该方法将深度学习和有限元结合起来，在保留一定的精度的同时实现了效率的极大提高，为虚拟手术模拟提供了一种高效的建模方法\cite{ref1}。刘代清从手法淋巴引流的角度创建了皮肤软组织力学模型，并用有限元做了淋巴动力学的分析，拓展了生物力学建模在康复和医疗方面的应用范围\cite{ref2}，张天明对人类运动进行肌骨系统生物力学建模和仿真研究，给运动分析、康复工程提供力学依据\cite{ref3}。吕晓庆等人对人体生物力学仿真和可视化相关研究做了全面的梳理，总结出多体建模、有限元、肌骨仿真等主要的技术路线，为人体交互仿真的方法提供参考\cite{ref4}，陈卓利用OpenSim平台概述了人体肌骨仿真建模的发展历程，使仿真工具在体育科学、康复工程领域向着标准化的方向发展\cite{ref5}。

潘言心设计并制作出分布式三维力触觉传感器，解决了分布式接触力精确采集的问题，给高分辨率触觉交互提供硬件支持\cite{ref6}。陈大鹏等人提出依靠数据的纹理摩擦建模和触觉渲染方法，通过采集真实物体表面摩擦特性来提高触觉渲染的真实性、材质区分度\cite{ref7}，陈庚对虚拟纹理全面属性建模和力触觉重现展开研究，将触觉反馈从单一的力值扩展为粗糙度、摩擦、纹理等全面特征，再次提升了虚拟交互的沉浸感\cite{ref8}，王宏润关注XR自然手势交互，创建起触觉加强系统，探寻手势识别和触觉反馈融合的方法，改进沉浸式交互的自然程度\cite{ref9}。

刘浩城针对针灸机器人软组织安全交互难点，提出柔性末端控制方法并重视力觉约束和安全策略，给医疗场景下柔性人机交互提供安全范例\cite{ref10}，梅欣根据物体力触觉属性识别，对机器人自适应抓取进行研究，用触觉感知提高复杂物体操作的稳定性、鲁棒性\cite{ref11}。

马雨前把RPIM无网格法应用到虚拟手术软组织形变仿真中，改进了碰撞检测算法，既保证了仿真精度又保证了实时性\cite{ref12}，冯上涛从虚拟手术缝合过程出发，加强了软组织形变仿真和碰撞检测的真实性、实时性，更符合临床手术操作的要求\cite{ref13}，吕双祺等对气凝胶复合材料做细观力学建模和变形行为仿真，扩大了可变形物体建模在先进功能材料领域中的应用范围\cite{ref14}，李浩达对多体系统非理想运动副进行接触力学建模与仿真分析，完善了复杂约束下动力学求解的方法\cite{ref15}，赵飞对重载履带式爬壁机器人进行力学建模与仿真，为大型装备柔性接触操作提供动力学支持\cite{ref16}。

国内的研究在力学建模的改良、触觉渲染维度的拓展、医疗和机器人场景的应用、硬件传感器的设计等各方面都取得了很多成果，但是系统级视触同步、多线程实时架构、跨设备标准化融合、全流程定量评价体系以及完整的``建模-交互-评价''循环平台创建还存在很大的改进空间。

\subsection{国外研究现状}

国外对于可变形物体触觉交互、实时仿真、数字孪生、柔性机器人、多模态交互等已经有了较早的研究，已经形成了一套包含基础理论、算法框架、系统整合和量化评定的技术体系，非常重视实时性、鲁棒性和工程应用能力。

Tzimos
N等对虚拟现实触觉接口和多模态交互进行研究，探究VR环境下3D物体的触觉接口多模态交互机制，证明触觉反馈可以明显提高交互的可靠性和准确性，也可以提高用户沉浸感\cite{ref17}。Taghizad
H等人的实时仿真理论及框架综述了多速率实时仿真技术、模型、框架以及主要难点，为高保真、高实时性交互系统规划提供理论依据\cite{ref18}。

Laukaitis
A等人把Webots与MuJoCo结合使用生成式AI来改善仿真的机器人学习和智能控制的实用价值，给多引擎仿真协同开拓了新途径\cite{ref19}，Islam
U N M
等人把人体姿态稳定和神经控制结合起来，用键合图模型提高姿态稳定性预测精度\cite{ref20}。

航天和极端环境仿真中，Han Z Z
等人用变分积分器建立小行星探测器的轨道动力学模型并进行仿真，从而提高长时间仿真的数值稳定性\cite{ref21}。Mao
Z等人站在系统动力学的角度研究教师数字素养和革新的影响机理，给复杂的社会系统仿真建模提供理论上的借鉴\cite{ref22}，针对复杂的群体运动仿真，Zhang
B等人给出滑雪场景下行人动力学建模及仿真的框架，给复杂的运动交互提供可扩展的方法\cite{ref23}。

柔性机器人建模与控制研究中，柔体外翻式生长机器人动态建模、物理分析及仿真实验验证由Kalibala
A等人完成，使软体机器人可以得到准确的控制\cite{ref24}。为了解决精密设备在振动环境下产生抖振等问题而建立的空间磁悬浮隔振控制系统进行动力学建模和仿真研究，Ye
M等人提出了一种新的提高精密设备稳定性的方法\cite{ref25}。对于可变形物体的形状控制，Zueco
C
I等人给出了一个多尺度弹性轮廓映射的方法，可以实现对复杂可变形物体高精度的形状调整和跟踪\cite{ref26}。

国外研究大多集中在高实时性保障、多速率同步机制、严格量化评价体系、跨平台硬件融合等几个方面，许多工作直接采用``物理引擎+触觉设备+VR''循环结构，表现出当下可变形物体触觉交互仿真领域的主要技术走向。

\subsection{研究述评}

综合国内外的相关研究可知，可变形物体触觉交互、实时仿真在力学建模、触觉渲染、多模态融合、柔性机器人、医疗仿真等各方面已经取得了很多成果，但是目前的研究成果还不能完全满足沉浸式、高真实感、低时延的实际交互需求。

可变形物体具有非线性，大形变的特性，实时性与物理保真度很难同时满足，轻量化模型会降低触觉真实感，高精度模型又不能达到触觉交互需要的高刷新率和低时延。物理引擎接触信息、触觉设备缺少对可变形物体的专门映射逻辑，力反馈、纹理渲染大多采用简单的线性映射，材质区分度、触感自然度与真实的物理交互存在差异。

视觉渲染与触觉反馈的时间没有共同合作，在时间上不同步，使得各个模块有不同的速度，从而产生延迟、抖动、体感错位等现象，从而影响到操作的稳定性以及沉浸感。

大部分研究只对单个模块或者视觉效果的表现形式进行关注，很少有人会制定统一的定量评价标准并做对照实验，对于任务效率、力量稳定性、系统时延、材质识别准确率等重要数据并没有进行系统的收集和验证，因此研究成果很难被重现，也无法进行横向比较。

目前该领域缺少创建模型规范、反馈真实、同步精准、性能稳定且支持全流程定量评定的可变形物体实时触摸仿真平台，本课题就从这里出发，用MuJoCo作为物理仿真核心、WEART作为触觉终端、VR作为视觉通道来构建视触融合的系统，通过多线程框架、自适应触觉渲染、改进后的逆运动学和标准化实验评价来弥补现有的技术缺陷，为虚拟手术、柔性操作、远程触觉交互等应用提供标准平台和技术参照。

\section{研究内容与技术路线}

\subsection{研究内容}

使用MuJoCo的flexcomp创建grid、ellipsoid、gmsh这三种可变形物体模型，并设置包含弹性、阻尼、质量、摩擦等参数化的配置，这些类型包含软体、布料、粘弹性组织等常见的种类。

触觉反馈机制的设计包括创建力值计算、纹理映射两个通道的触觉渲染功能，可以达到接触力归一化处理、自动适应性补偿、滑动过滤、动态纹理调节的目的，给用户带来低时延高质量的反馈响应。

运动重定向和逆运动学相结合，可以使VR姿态以及手指闭合度准确地映射到虚拟手上来，用阻尼最小二乘法（DLS）改进IK求解过程，从而提高稳定性，减少关节超出范围的情况。设计多线程并行结构，把仿真、触觉、渲染三线程分离且同步，从而达到高刷新率以及低端到端的延迟。

实验设计和量化评定，得到纯视觉/视触融合对照实验，用任务历时、达成率、力的稳定性、系统滞后、帧率等角度进行验证和分析。

\subsection{技术路线}

本研究遵循建模---交互---评估闭环技术路线：

(1)理论与文献梳理，整理出可变形建模、触觉渲染、实时仿真等技术的现状，确定技术方案；

(2)系统设计，确定硬件接口、软件模块、数据流和算法模块；

(3)实现和集成，完成可变形模型、触觉反馈、IK、多线程同步开发；

(4)测试与优化：功能测试、性能基准测试、参数调优；

(5)实验和分析阶段对用户进行对照实验，统计分析并检验触觉反馈的有效性；

(6)总结与展望：归纳成果、局限与未来方向。
